% Part: sets-functions-relations
% Chapter: arithmetization
% Section: reflections

\documentclass[../../../include/open-logic-section]{subfiles}

\begin{document}

\olfileid{sfr}{arith}{ref}
\olsection{Sawetara Pamawas Filosofis}

Wis cukup perkara teknise. Nanging apa kang wis digayuh?

Kang meh ora bisa dibantah, kita oleh matematika murni kang {endah}.
Kajaba iku, ana sawetara gegayuhan konseptual kang jero. Pangerten yen
Sipat Kajangkepan ngandharake prabédan pokok antarane wilangan real lan
wilangan rasional iku wawasan kang jero. Kajaba iku, pambangunan wilangan
real kanthi cetha minangka potongan Dedekind menehi dhasar kang kukuh
kanggo analisis. Kita ngerti yen gagasan \emph{lapangan katata kang
jangkep} iku koheren, amarga potongan-potongan mau mbentuk lapangan kaya iku.

Sanadyan mangkono, kita kudu ngandharake sawetara cathetan kritis tumrap
gegayuhan kasebut.

Kapisan, durung cetha yen mikirake wilangan real nganggo potongan iku
\emph{luwih} ketat tinimbang mikirake wilangan real nganggo ekspansi
desimal kang lumrah (lan bisa tanpa wates). ``Pambangunan'' wilangan real
kang kapindho iki rada mirip karo pambangunan wilangan real nganggo
rerangken Cauchy; nanging sejatine cara iku wis dimangerteni para
matematikawan wiwit awal abad kaping 17 (delengen \olref[cauchy]{sec}).
Tambahing kaketatan kang nyata dumadi nalika dimangerteni yen wilangan
real nduweni Sipat Kajangkepan; kabisan mbangun wilangan real minangka
himpunan tartamtu mbokmenawa ora pati narik kawigaten yen ngadeg dhewe.

Luwih ora cetha maneh yen aritmetisasi wilangan wutuh utawa wilangan
rasional, kang luwih gampang, nambah kaketatan ing babagan kasebut.
Ing kene ana pangamatan prasaja kang prayoga. Sawise \emph{mbangun}
wilangan wutuh minangka kelas ekuivalensi pasangan urut wilangan natural,
banjur mbangun wilangan rasional minangka kelas ekuivalensi pasangan urut
wilangan wutuh, lan banjur mbangun wilangan real minangka himpunan wilangan
rasional, kita langsung \emph{nglalekake} pambangunan kasebut. Mligine,
ora ana wong kang kepengin \emph{nggunakake} pambangunan iki sajrone bukti
matematis (kejaba, mesthi wae, bukti yen pambangunane tumindak kaya kang
kudune). Luwih gampang ngomong langsung bab wilangan real tinimbang bab
sawijining himpunan saka himpunan saka himpunan saka himpunan saka himpunan
saka himpunan saka himpunan wilangan natural.
%(Amarga wilangan real dibangun minangka himpunan wilangan rasional; lan wilangan rasional dibangun minangka kelas ekuivalensi pasangan urut wilangan wutuh, yaiku himpunan saka himpunan saka himpunan wilangan wutuh; lan sateruse.)
%2/3 = [2,3]
%	yaiku {<2,3>, ... }
%	{ {{2},{2,3}}, ... }
%
%wilangan real iku
%	himpunan wilangan rasional
%
%wilangan rasional iku
%	himpunan saka himpunan saka himpunan wilangan wutuh
%
%wilangan wutuh iku
%	himpunan saka himpunan saka himpunan wilangan natural

Kang paling mamang yaiku apa definisi-definisi iki ngandhani kita apa
\emph{sejatine} wilangan wutuh, rasional, utawa real, \emph{yen dirembug
saka sisi metafisika}. Tegese, kita patut mangu yen wilangan real, umpamane,
\emph{iku} himpunan tartamtu (saka himpunan saka himpunan\ldots). Pepalang
utama tumrap panemu iki yaiku pambangunane bisa ditindakake kanthi maneka
warna cara. Tumrap wilangan real, ana sawetara pambangunan beda kang pancen
narik kawigaten (delengen \olref[cauchy]{sec}). Nanging iki cara kang prasaja
banget kanggo ngolehake pambangunan beda: kaya ing
\olref[sfr][rel][ref]{sec}, kita bisa netepake pasangan urut kanthi cara
rada beda; yen gagasan pasangan urut alternatif iki digunakake,
pambangunan kita tetep bisa lumaku persis padha becike, nanging pungkasane
ngasilake obyek kang beda. Mula ana akeh pambangunan teori-himpunan kang
saingan kanggo wilangan wutuh, rasional, lan real. Saiki bakal sawenang-wenang
(lan gawe isin) yen kita ngaku yen wilangan wutuh, umpamane, yaiku
\emph{himpunan-himpunan iki}, dudu \emph{himpunan-himpunan kuwi}. (Kaya ing
\olref[sfr][rel][ref]{sec}, iki tuladha argumen kang digawe misuwur dening
\citealt{Benacerraf1965}.)

Ana perkara liyane kang prayoga dirembug: pambangunan kita nduweni bab kang
cukup \emph{aneh}. Kita miwiti saka wilangan natural. Banjur kita mbangun
wilangan wutuh, lan mbangun ``$0$ saka wilangan wutuh'', yaiku
$ \equivrep{0,0}{\Intequiv}$. Nanging $0 \neq
\equivrep{0,0}{\Intequiv}$. Pancen, miturut pambangunan kita,
\emph{ora ana} wilangan natural kang dadi wilangan wutuh. Nanging iki katon
ora lumrah banget. Pancen, ing \olref[sfr][set][imp]{sec}, kita ngaku tanpa
akeh argumen yen $\Nat \subseteq \Rat$. Yen pambangunan kasebut ngandhani
kita persis \emph{apa} wilangan-wilangan iku, pratelan iki mesthi kleru.

Yen saiki kita mundur sedhela, tekan ngendi asile? Kanthi makarya ing teori
himpunan naif lan nganggep wilangan natural wis cumawis, kita bisa
\emph{nganggep} wilangan wutuh, rasional, lan real minangka himpunan-himpunan
tartamtu. Ing teges iki, kita bisa \emph{nyisipake} teori bab obyek-obyek iki
ing sawijining teori himpunan. Nanging wigati filosofis saka panyisipan iki
ora prasaja.

Mesthi wae, iki kabeh dudu tembung pungkasan!{} Pokoke mung mangkene.
Nuduhake yen aritmetisasi wilangan real \emph{pancen} wigati banget saka
sisi filsafat mbutuhake argumen \emph{filosofis} tambahan.

%Kajaba iku: sajrone bab iki, kita nganggep wilangan natural mung wis
%\emph{diwenehake} marang kita. Nanging apa kang bisa ditindakake
%nganggo wilangan kasebut? Iku bahan kanggo bab sabanjure.

\end{document}
